\documentclass[
  a4paper,
  12pt,
  oneside,
  parskip=half,
  listof=totoc,
  bibliography=totoc,
  DIV=12,
  numbers=noenddot,
  BA.tex
  ]{subfiles}

\begin{document}

\section{ALICE at the Large Hadron Collider}

The Large Hadron Collider (LHC) is the world's largest and most powerful particle accelerator, installed in a $26.7\bunit{km}$ underground tunnel at the CERN research complex near Geneva, Switzerland. \cite{bruning_large_2012} The particles are accelerated in two beam-pipes which can cross in four interaction regions, allowing particles of the same charge, for example proton-proton (pp) or lead-lead (Pb--Pb) to be collided. The LHC is capable of accelerating particles to a center-of-mass energy of up to $14\bunit{TeV}$ while maintaining a luminosity of order $10^{34} \frac{1}{\bunit{cm}^2 \cdot s}$ for pp collisions. The LHC ring is built with superconducting NbTi magnets, operated at a superfluid helium temperature of $1.9\bunit{K}$, allowing magnetic fields of up to $B = 8.33\bunit{T}$ for keeping the accelerated particles on track. The main goal of the LHC is to probe the Standard Model and physics beyond the Standard Model (BSM) in higher energy regimes than ever done before using hadrons, like protons, as collision particles. The LHC is also capable of colliding heavy ions, like stripped lead ions ($^{208}\text{Pb}^{82+}$), which are principally provided at Interaction Point 2 for the specialized ALICE detector. \cite{evans_lhc_2008}

\begin{figure}[H]
	\centering
	\includegraphics[width=0.9\textwidth]{Media/Pictures/ALICE_RUN3_no_labels_no_ITS.jpg}
	\caption{The ALICE experiment at the LHC, CERN during RUN 3. Image taken from \cite{cern_document_server_alice_nodate}.}
	\label{fig:ALICE_CAD_Overview}
\end{figure}

ALICE (A Large Ion Collider Experiment) is one of the biggest experiments at the Large Hadron Collider. \cite{collaboration_alice_2008} It is a dedicated heavy-ion experiment with the focus on studying QCD matter and especially the QGP. ALICE also serves other purposes than taking data in Pb--Pb ion collisions. Important tasks are collisions with lighter ions, proton-nucleus runs and pp beam data taking as a reference data measurement for heavy ion physics. Lighter ion collisions include oxygen-oxygen collisions that were first conducted in 2025. \cite{alice_collaboration_oxygen_2025}

The ALICE detector itself has a diameter of 16m, a length of 26m and weighs about 10.000 tons. \cite{collaboration_alice_2008} It provides a high spatial resolution and a low transverse momentum threshold of $p_T^{min} \approx 0.15 \bunit{GeV}$, meaning particles with low transverse momenta produced in the collision can be detected reliably. Furthermore, it serves a good particle identification for particles with transverse momenta of up to $20 \bunit{GeV}$. ALICE is composed of 18 detector systems (see figure \ref{fig:ALICE_CAD_Overview}) while each one serves a very specific design concept purpose for the experimental conditions expected at the LHC. One major criterion for the ALICE detector systems is to cope with the extreme particle multiplicity occurring in central Pb--Pb collisions. ALICE consists of the central barrel part, used for measuring hadrons, electrons and photons and a forward muon spectrometer. The central part houses the Inner Tracking System (ITS) made up of high-resolution silicon pixel detectors, the cylindrical Time Projection Chamber (TPC), three Time-of-Flight (TOF) sensor arrays, Ring Imaging Cherenkov (HMPID) detectors, Transition Radiation (TRD) detectors and two electromagnetic calorimeters, all embedded in a large solenoid magnet with a magnetic field of $B = 0.5\bunit{T}$. The magnetic field is necessary because the trajectories of charged hadrons are bent to an arc in the magnetic field. By measuring the radius of the arc it is possible to measure the ratio between the charge and the momentum of the hadron, which is needed for particle identification.

Since collisions in ALICE have such a high multiplicity, the most challenging task is reconstructing the particle tracks. \cite{collaboration_alice_2008} To do so, the primary vertex of the collision is reconstructed using the hit positions in the innermost silicon pixel detector (SPD) firstly. Then, the primary vertex's position is used to initiate the tracking of primary particles in the TPC and for the subsequent track following. Secondary tracks and vertices are found during a further tracking pass. In general, charged hadrons are identified using the information of the ITS, TPC, TRD, TOF and HMPID detectors, while electrons are identified using the TRD, TPC and EMCal detectors. Photons are detected in the PHOS detector, as well as neutral particles, like $\pi^0$ or $\eta$. Neutral particles are detected by taking $\gamma \gamma$ invariant mass spectra. Muons are detected in the forward muon spectrometer at at pseudo-rapidity region of $-4.0 < \eta < -2.5$. 

\end{document}